% ---
% id: reductions/wpfss/prop:weight-monotone
% kind: prop
% title: "Monotonicity of the weight-preserving target in its weights"
% summary: >-
%   Certified weights stay certified when lowered but not below zero, and while
%   the total weight is below one the weight-preserving target of the larger
%   weights squashes to that of the smaller ones. By [[cor:monotone]] the
%   smaller weights give the cleaner target, which hands the adversary more.
% prominence: standard
% status: accepted
% uses: [def:targets, prop:chain, reductions/cor:monotone]
% mentions: []
% assumes: []
% refs: []
% provenance: newly-derived
% lean:
% ---

Nothing in Definition~\ref{def:targets} or Proposition~\ref{prop:chain}
requires $p_v$ to be the largest admissible value. The proposition says what
lowering the weights does: they stay valid, and the target becomes cleaner, a
statement about the pre-order that is proved rather than asserted.

\begin{proposition}
\label{prop:weight-monotone}
Let $(p_v)_{v\in\Vset}$ be \tref{certified weights} for $M$ with
\tref{total weight} $p$, let $p'_v\in[0,p_v]$ for every $v\in\Vset$, and
write $p'\coloneqq\sum_vp'_v$. Write $\Fwp[p]$ and $\Fwp[p']$ for the
\tref[weight-preserving target]{weight-preserving targets} \eqref{eq:WPFSS}
built from $(p_v)$ and from $(p'_v)$.
\begin{enumerate}[label=(\roman*),leftmargin=3.0em]
\item $(p'_v)$ are certified weights for $M$, and $p'\le p$; in particular,
  where $p<1$ the construction of Proposition~\ref{prop:chain} runs
  unchanged with $(p'_v)$.
\item If $p<1$, put
  \begin{equation}\label{eq:weightmono}
    \Psi(\sigma)\coloneqq\lambda\,\sigma+(1-\lambda)\,\Pleq\sigma\Pleq
     +\tr\!\big[\Pflag\sigma\big]\,\nu,
    \qquad
    \lambda\coloneqq\frac{1-p}{1-p'},
    \qquad
    \nu\coloneqq\sum_v\frac{p_v-p'_v}{1-p'}\dyad{v}_{\Flag}.
  \end{equation}
  Then $\Psi$ is a channel and $\Fwp[p]\dsq_{\Psi}\Fwp[p']$, so by
  Corollary~\ref{cor:monotone} the smaller-weight target is the cleaner one.
\end{enumerate}
\end{proposition}

\begin{proof}
(i) By \eqref{eq:certweights},
$\blk{M}{\gtN}_v\succeq p_v\idon{\Bsys_{\gtN}}\succeq p'_v\idon{\Bsys_{\gtN}}$
with $p'_v\ge0$, for every $v$, and summing $p'_v\le p_v$ gives $p'\le p$; so
where $p<1$ also $p'<1$, and Proposition~\ref{prop:chain} applies to
$(p'_v)$.

(ii) \emph{$\Psi$ is a channel.} $\lambda\in[0,1]$ and $\nu\succeq0$ because
$p'\le p<1$ and $p'_v\le p_v$, so $\Psi$ is CP; and with
$\tr\nu=(p-p')/(1-p')$ and $\Pleq+\Pflag=\idon{\Qsys}$,
$\tr\Psi(\sigma)=\tr[\Pleq\sigma]+\big(\lambda+\tr\nu\big)\tr[\Pflag\sigma]
=\tr\sigma$, since
$\lambda+\tr\nu=\big((1-p)+(p-p')\big)/(1-p')=1$.

\emph{The adjoint.} On block observables
$\Psi^{\dagger}(X_{\leqN}\dsum X_{\Flag})
=X_{\leqN}\dsum\big(\lambda X_{\Flag}+\tr[\nu X_{\Flag}]\idon{\Flag}\big)$, and
evaluating at $X_{\Flag}=p'_v\idon{\Flag}+(1-p')\dyad{v}$ gives flag block
\begin{equation*}
  \Big[\tfrac{(1-p)p'_v}{1-p'}+\tfrac{p'_v(p-p')}{1-p'}+(p_v-p'_v)\Big]\idon{\Flag}
  +(1-p)\dyad{v}
  \;=\;p_v\,\idon{\Flag}+(1-p)\dyad{v},
\end{equation*}
because the first two brackets sum to $p'_v$; with the unchanged $\leqN$ block
this is $\Fwp[p]_v$. The conclusion is Corollary~\ref{cor:monotone} applied
to $\Fwp[p]\dsq_{\Psi}\Fwp[p']$.

Taking $p'\equiv0$ recovers the flag-noise channel $\Lnoise$ of
\eqref{eq:Lnoise}, as it must.
\end{proof}
